\documentclass[]{../../tex/classes/styledArticle}
\usepackage{../../tex/packages/hebrewSupport}
\usepackage{../../tex/packages/mathShortcuts}
\usepackage{../../tex/packages/theoremsSupport}
\usepackage{listings}

\author{שחר פרץ}
\title{פרויקטו 1}
\date{16 באפריל 2026}
\begin{document}
	\maketitle
	יש שני תרגילי בית, כל אחד בהיקף 5\%, פרויקט 45\%, ומבחן 45\%. ההרצאות ביום רביעי יתקיימו במתכונת של כיתה הפוכה – מתחרים בזום לשאול שאלות על ההרצאה מיום חמישי. תרגילי בית 0 לא נכנס לציון. 
	
	עם השנים, רואים ירידה בהתמודדות העצמית של הסטודנטים עם הקוד, וזה משתקף במבחן. זה לא רק לכתוב קוד עם סינטקס בעצמכם, זו גם היכולת לגשת ולפתור את השאלות. זה בין הקורסים המרכזיים בתואר של תכנות, ההמלצה של המרצה היא לפתור את תרגילי הבית בעצמכם. כלי AI נהיים הרבה יותר טובים כשאתם יודעים לתכנת בעצמכם ומבינים את הפלט שלהם. רואים את ההבדל בין סטודנט שמבין את החומר לבין סטודנט שלמד כל הסמסטר עם chatGPT. 
	
	בשיעורי הבית הבדיקות הן אוטומאטיות. אם שם הקובץ לא נכון – בום כל הנקודות יורדות וקיבלתם 0. כנ''ל בנוגע לשגיאות קומפילציה. יש מנגנון ערעורים בקורס, שב־99\% זה לא ''הבודק טעה`` אלא תיקוני שורה או שתיים שיכולות להביא אותכם מ־0 ל־100, עם קנס של 5-10 נקודות לתיקון. לכל קבוצה יש ערעור אחד. תוודאו שהערעור מתקמפל, וכמובן שהוא מפורט. 
	
	היום נכיר את סביבת העבודה וכן python בסיסי. מחברות ה־Jupyter נמצאות באתר. 
	
	אם רוצים לעשות cp ממקומי לשרת מרוחק שדורש סיסמה, יש להשתמש ב־scp: 
	
	
	\sen\begin{lstlisting}
scp <local\_file> <usename>@nova.cs.tau.ac.il:/home
	\end{lstlisting}\she
	
	\sen
	\subsection*{Docker}
	\she
	
	בניגוד ל־VM, תהליך מאוד כבד  ומסורבל, Docker עובבד בשיטה של containers. הוא מדמה סביבה של מערכת הפעלה נפרדת, על אף שהוא לא מעלה מערכת הפעלה שלמה. הוא מאשר לנו ליישר קו בין סביבות עבודה. אנחנו לא רוצים לפול על זה שהשתמשנו בקומפיילר שונה משל הבודקים, או שחסרה חבילה, או משהו כזה. חשוב מאוד לעשות zip על ה־docker, זה משתנה פר סביבת עבודה. הסיבה? על Mac ו־Windows יש כל מני חתימות ב־zip שלא מופיעות לבודקים האוטומאטיים. 
	
	מה זה בכלל אומר מחשב מרוחק (למרות שהוא רץ מקומית, הוא מתחבר ל־socket על הרשת הלוקאלית של המחשב)? ה־Docker מתחיל משורת פקודות שנקראת Dockerfile. ה־Dockerfile של הקורס נמצא במודל. הוא מגדיר ומריץ כל מני חרא ברגע שהוא עולה. 
	
	נרצה ליצור image. אפשר לחשוב על זה כ''קובץ הרצה``. התוכנית תקומפל ל־image לפי ה־dockerfile. איך עושים את זה? באמצעות: 
	\sen\begin{lstlisting}
docker build -t c-python:latest .
	\end{lstlisting}\she
	כאשר latest תמיד הגרסה האחרונה, ובמקומה אפשר לשים גרסה ספציפית. ה־. משמעה התיקייה הנוכחית, כלומר הוא מחפש את ה־image בתיקייה הנוכחית.
	
	וודאו שה־dockerfile בלי סיומת של .txt, למרות שככה הוא יורד מגיטהאב. 
	
	אחרי כן יש להעלות את הקונטיינר באמצעות docker run. המשמעות של -d שהוא רץ ברגע, והמשמעות של -p הוא למפות port מהמחשב ל־port ב־docker. המשמעות של -v היא volume – מיפוי מתיקייה (Volume) המחשב הנוכחי לתיקייה של הדוקר. לבסוף יש לציין איזה image בכלל אנו רוצים להריץ. 
	
	אנחנו ''תכלס`` צריכים לעשות docker run רק פעם אחת בקורס. תמיד אפשר להריץ על ה־Nova, במידה ויש לכם מחשב ישן מדי ויש בעיה ב־docker לא פתירה. ואם מישהו נתקע, יש שעות קבלה במודל. היום ה־Nova היא אופציונלית בלבד. 
	
	
	\ndoc
\end{document}